\chapter{Literature Review and Theoretical Foundation}

%% English draft. \crosscheck{...} flags content requiring human verification.
%% All \cite{...} keys resolve to references.bib.
%%
%% The Related Work section below was built from a literature search. Every claim about
%% another author's system has since been checked against that source's full text, and the
%% \fromabstract{} daggers were discharged accordingly. Three claims did not survive the
%% check and were rewritten: ke2025 (the hallucination-free run was GPT-4 over
%% international guidelines alone, not all ten models), degroot2024 (the "real EHR data"
%% sample is one patient) and kraljevic2021 (the 8.8B tokens are one hospital's
%% self-supervised training corpus, not a three-hospital validation set). Two \crosscheck
%% blocks (Sections 2.1.2 and 2.1.4) still flag findings that QUALIFY the novelty claim
%% made in Chapters I and IV; zhou2022 and ali2026 both hold up exactly as described, so
%% those two chapters still need revising before submission.

\section{Related Work}

This section reviews prior work by theme and locates the gap this thesis addresses.

\subsection{Guideline-based clinical decision support}
Computerised decision support is an established route to narrowing the gap between
guidelines and practice, and its effect size is now well quantified. A meta-analysis \cite{kwan2020} of 122 trials covering 1{,}203{,}053 patients found that
decision support increased the proportion of patients receiving recommended care by 5.8\%
(95\%~CI 4.0--7.6), with substantial heterogeneity ($I^2 = 76\%$). Improvements in clinical
endpoints were far smaller, at a median of only 0.3\%. The value of such systems is therefore real but modest, and concentrated in process rather
than outcome measures. The same pattern appears in perioperative settings, where adherence
improved and some safety outcomes improved with it, but postoperative mortality did not
\cite{cai2024}.

The dominant implementation tradition is that of \emph{computer-in\-ter\-pret\-able guidelines}
(CIGs), in which guideline logic is encoded in a dedicated formalism---Arden Syntax,
GLIF, PROforma, Asbru, EON and later GLARE and SAGE---and executed by an engine against
patient data \cite{declercq2004,declercq2008,oliveira2014}. No single formalism has been
adopted as a standard \cite{declercq2008}, and Oliveira \textit{et al.}
identify the interoperability between a guideline execution engine and the clinical data
model as a persistent obstacle \cite{oliveira2014}---precisely the seam
this thesis addresses by matching guideline concepts to coded EHR data through a shared
concept space.

Empirically, that seam has remained open. A systematic review of rule-based decision
support screened 1{,}731 papers and found only 15 reporting
any evaluation of clinical outcomes; among those, only two integrated with
an EHR system, six used ontologies, and three used a recognised healthcare technology
standard \cite{papadopoulos2022}. Reported barriers across this literature
are consistent: opacity of encoded rules, alert fatigue, poor workflow fit, and weak
evaluation methodology \cite{papadopoulos2022,bosetti2025,silva2022}.
Two observations follow for this thesis. First, EHR integration is the exception rather
than the rule, which is what makes a FHIR-native design a contribution rather than an
implementation detail. Second, the scarcity of rigorous evaluation makes the
concordance-against-an-independent-reference design of Chapter~III a methodological
contribution in its own right.

\subsection{Knowledge graphs for clinical decision support}
\label{sec:rw-kg}
Knowledge graphs integrate heterogeneous clinical knowledge and patient data into a
queryable structure, and the healthcare literature on them is now substantial enough to
have been reviewed several times \cite{abusalih2023,cui_h2023,budhdeo2023}. A scoping
review of 255 articles reports median graph sizes of
46{,}983 nodes (IQR 6{,}415--460{,}948) and
906{,}737 edges (IQR 66{,}272--9{,}894{,}909), and concludes
that translation of knowledge graphs into clinical practice remains a
challenge \cite{budhdeo2023}. The graph reported in Chapter~IV sits above that median on
both counts, in the upper part of the reported interquartile range.

Prior systems divide along an axis of what they connect. On one side are \emph{patient}
graphs: Xiao \textit{et al.} expose an OMOP repository as a FHIR RDF knowledge graph and
validate it by query equivalence against the source database, but
encode no clinical guideline and therefore issue no recommendations
\cite{xiao2022}; Shang \textit{et al.} build an EHR-oriented graph for collaborative
reasoning across institutions, applied to identifying overlooked chronic kidney disease
\cite{shang2024}. On the other are \emph{knowledge} graphs: Cui
\textit{et al.} build a cardiovascular medication guideline graph in Neo4j but
attach no patient records \cite{cui2025}. Earlier FHIR-based decision aids
such as that of Semenov \textit{et al.} pair a FHIR interoperability layer with a
graph-structured knowledge base, but operate as production-rule systems over
a narrow laboratory panel rather than over longitudinal records \cite{semenov2018}.
Seneviratne \textit{et al.} are closer in spirit to the present work, modelling guideline
provenance and study cohorts so that recommendations can carry their evidence
\cite{seneviratne2023}.

A more mundane body of work must also be named here, because it bounds what this thesis may
claim. The mapping from FHIR to a property graph---resources become nodes, references become
edges---is mechanical and well understood, and exists as commodity tooling: browser-based
bundle viewers render an uploaded bundle as an interactive node--edge graph without asking any
modelling decision of the user \cite{tenasol_viewer}. The transformation itself is therefore
not claimed as a contribution of this work. What such tools do \emph{not} do marks the
boundary. They operate on one bundle at a time and do not persist it, so no cross-patient
question can be posed. They render each coding inline within the resource carrying it rather
than hoisting codings into shared nodes, so no two patients are joined through a common code.
They reconcile nothing between SNOMED~CT, RxNorm and LOINC. And they contain no knowledge that
did not arrive inside the bundle, so there is no guideline or literature layer to reason
against. A viewer is moreover required to be \emph{faithful} to its input, inspection being its
purpose. The graph built in Chapter~III, by contrast, is deliberately lossy: it discards
billing and provenance resources and all but nine LOINC observation codes. The reason is
that it is built to answer one clinical question rather than to represent a document.

One published system spans both sides of that axis, and it bounds what this thesis may
claim. Zhou \textit{et al.} construct a hypertension-medication knowledge graph from the
Chinese hypertension guideline and evaluate it against 124 health records
collected from a single chief physician at one hospital, reporting recall
91\%, hit@3 83\% and MRR 77\% on a single-drug-class recommendation task \cite{zhou2022}.
A guideline layer, patient records and a recommendation evaluation therefore already
coexist in one system, and that combination is \emph{not} claimed as novel here. What
separates the present work is scope and evaluation rather than architecture, and the
margin is narrower than their abstract suggests. Their graph spans every antihypertensive
drug class and carries the guideline's compelling level into its scoring rules. Neither
multi-class coverage nor recommendation strength therefore distinguishes the present work.
The breadth of the guideline encoded does: it spans antiplatelet, lipid, blood-pressure,
beta-blocker, renin--angiotensin--aldosterone, lifestyle and contraindication-safety
recommendations rather than drug therapy for a single condition. Numeric physiological
values likewise enter their reasoning path: physical attributes are one of the four entity
types in their graph, and an indication may be qualified by a threshold, as in their own
example of a dihydropyridine calcium-channel blocker indicated for stable angina and
hypertension at a heart rate below 75~bpm. What separates the two is how those values
arrive---keyed by a clinician into a patient-entry form for 124 records there, and here as
FHIR \emph{resource} instances carrying their own units and code systems, at cohort scale
and without manual transcription. Their evaluation reports ranking metrics over a candidate
list, whereas concordance here is measured against a CRS-independent reference
standard and a manually labelled gold standard. Finally, their system fuses two sources; a
literature layer is fused here as a third.

\subsection{Terminology integration}
Matching coded patient data to guideline concepts requires reconciling several
vocabularies---SNOMED~CT, RxNorm and LOINC---for which the UMLS Metathesaurus is the
standard bridging resource \cite{umls}. Mapping is not lossless: a comparison of
ATC-to-RxNorm mappings found 1{,}949 null mappings in UMLS against
916 in the OHDSI Standardized Vocabularies, though UMLS achieved higher
coverage (96.5\% versus 91\%) on a single-patient EHR sample, a difference that
traced to two multi-ingredient ATC classes \cite{degroot2024}. Any concept-bridge
design therefore carries a residual mapping-failure rate that should be measured
rather than assumed.

For linking free text to concepts, MedCAT provides UMLS and SNOMED~CT concept
extraction, self-supervised over approximately 8.8 billion tokens of one hospital's
records and validated at three London hospitals, with reported transferability
$F_1 > 0.94$ between sites \cite{kraljevic2021}. This work
uses scispaCy \cite{scispacy} instead, for reasons of task fit rather than of quality.
MedCAT's advantage is contextual disambiguation learned self-supervised from a corpus of
clinical narrative. Neither condition holds here. The bridging task of
Section~\ref{sec:kg-construction} runs over isolated code display strings---\texttt{Chronic
ischemic heart disease (disorder)}---which carry no surrounding sentence for such a model to
read. The FHIR export is moreover coded rather than narrative, so there is no corpus on
which to train one. The scispaCy library also exposes its knowledge base as a queryable alias index rather than
only as an annotator over text. That is what permits the exact-alias pass and the
semantic-type filter used here. Nor does the library hold any learned state, so a given
string yields the same concept on every run. This thesis depends on that property when it
claims that retrieval is deterministic. For free-text clinical notes, where abbreviations must be resolved from
their context, MedCAT would be the better instrument.

\subsection{Large language models, retrieval-augmented generation, and faithfulness}
Large language models generate fluent clinical text but hallucinate; retrieval-augmented
generation (RAG) mitigates this by conditioning generation on retrieved evidence
\cite{rag_lewis2020}. A systematic review of RAG in healthcare finds a range of
architectures in use but reports a lack of standardised evaluation
frameworks for RAG-based applications and that most studies do not address ethical
considerations \cite{amugongo2025}. That absence is what makes the
mechanical faithfulness check of Chapter~III worth reporting as a contribution.

Guideline-specific RAG systems now report strong results. A RAG system over the UK NICE
guidelines raised faithfulness---whether an answer is supported by the
source text---by 64.7 percentage points to 99.5\%
\cite{lewis2025}. A framework for hepatological guideline interpretation
improved accuracy from 43\% to 99\% through guideline reformatting and
prompt engineering \cite{kresevic2024}. An LLM-RAG evaluation across
ten models and 3{,}234 responses reported no hallucinations for its best
configuration, GPT-4 over international guidelines, which also outperformed
human-generated answers; hallucination rates across the other models ranged up to
48.6\% \cite{ke2025}.

Related work also pairs knowledge graphs with language models for explainability
\cite{park2025}, so the graph-plus-LLM combination is an active rather than an
unexplored area.

\subsection{Synthetic clinical data and FHIR-based research}
Synthea simulates patient lifespans through modular disease models and exports HL7 FHIR
bundles, making reproducible clinical-informatics research possible without privacy risk \cite{walonoski2017}. It has since been extended with imaging, genomic and note data
linked by FHIR \cite{walonoski2022}. Its validity has been examined
directly: comparing clinical quality measures computed on a 1.2-million-patient
Massachusetts cohort against publicly reported rates, Chen \textit{et al.} conclude that
Synthea is reliable in modelling demographics and the probability of
services being offered, but that its capabilities to model heterogeneous
health outcomes post services are limited and that it does not model
deviations in care \cite{chen2019}. Notably, they report that
no Synthea residents had their blood pressure controlled after a
hypertension diagnosis, against a Massachusetts rate of 74.5\%
\cite{chen2019}.

\subsection{Terminology and framing over time}
Searching this literature requires two vocabularies. Work published before roughly 2015
describes the same architecture in Semantic Web terms: OWL ontologies, description-logic
reasoners, RDF triples and proof engines. Examples are the ontology-driven preoperative
decision support of Bouamrane \textit{et al.} \cite{bouamrane2011} and the
guideline-execution framework of Hussain \textit{et al.} \cite{hussain2007}. The latter
already combined a domain ontology, a guideline ontology and a patient ontology, and
generated justification trees to expose the inference trace. Post-2020 work describes the same problem
using property-graph vocabulary---knowledge graph, Neo4j, Cypher, GraphRAG. The
architecture proposed in this thesis is therefore a modern-idiom instance of a
well-established design, and a search restricted to ``knowledge graph'' will silently miss
a decade of directly relevant precedent.

\subsection{Summary and research gap}
Prior work addresses these themes largely in isolation. Rule-based and CIG systems encode
guideline logic but seldom integrate with a standards-based EHR
\cite{papadopoulos2022}. Patient-side clinical graphs are rigorously
constructed but carry no guideline layer \cite{xiao2022}. Guideline-side graphs carry no
patient records \cite{cui2025}. One system does carry both, evaluating a guideline-derived
medication graph against real records \cite{zhou2022}, but for a single condition's drug
therapy and by ranking metrics rather than against a reference standard.
Graph-grounded language models demonstrably reduce
hallucination \cite{ali2026} but are evaluated on question-answering accuracy---sixty
questions in a single disease area, scored against reference answers drawn from the
literature---rather than on concordance with a guideline over a patient cohort. The contribution of this thesis is therefore narrower and more precise than ``joining
patient data to guidelines''. First, it builds a FHIR-native graph in which \emph{structured
observations} are available to gate recommendations. Second, that graph is evaluated for
guideline concordance against a CRS-independent reference standard. Third, it comes with a
grounded briefing whose citation-level faithfulness can be checked mechanically, because
retrieval is deterministic and provenance-bearing.


\section{Theoretical Foundation}

\subsection{Chronic coronary disease and the 2023 AHA/ACC guideline}
Chronic coronary disease (CCD) denotes the stable phase of coronary artery disease requiring
long-term secondary prevention.

\subsubsection*{A note on terminology: SIHD and CCD}
The two terms used in this thesis are related but not interchangeable, and the distinction
matters because the data and the guideline come from either side of it. Until 2023 the
condition was addressed under the heading \emph{stable ischemic heart disease} (SIHD),
denoting patients with established coronary disease who are not in an acute phase. The 2023
AHA/ACC guideline \cite{virani2023} replaced that framing with \emph{chronic coronary
disease}, a deliberately broader entity: alongside stable angina it explicitly encompasses
patients who are more than one year past an acute coronary syndrome, patients after
revascularization, patients with ischemic cardiomyopathy, and patients whose coronary
disease is documented without current symptoms. SIHD is therefore best read as a
\emph{subset} of CCD---every SIHD patient falls under CCD, but not every CCD patient would
have been called SIHD.

This asymmetry is visible in the present work. The synthetic corpus is generated with the
Synthea module named \texttt{stable\_ischemic\_heart\_disease}, whose vocabulary predates the
2023 renaming, whereas the encoded guideline is the CCD guideline. The module therefore
supplies patients drawn from the narrower, older category, while the recommendations are
written for the wider one. Two consequences follow. Cohort membership is not decided by which
generator module produced a patient but by the coded diagnosis actually present in the record
(Section~\ref{sec:gate}), so a patient from the general-population run who develops the
qualifying diagnosis is included on the same footing. And because CCD is the broader
category, recommendations addressed to parts of CCD outside the classical SIHD picture---for
example those keyed to a recent acute coronary syndrome---are exercised less often by this
corpus than they would be in a real CCD population.
The 2023 AHA/ACC CCD
guideline \cite{virani2023} organises its recommendations by \emph{Class of Recommendation}
(COR) and by \emph{Level of Evidence}. The COR is Class~1 (strong), Class~2a (moderate),
Class~2b (weak), or Class~3, the last being subdivided into 3-No-Benefit and 3-Harm,
i.e.\ contraindications. In this
work the COR determines the semantics of a recommendation edge: a Class-3 recommendation
becomes a contraindication, and Class~1/2a/2b recommendations become positive
recommendations.

\subsection{HL7 FHIR}
Fast Healthcare Interoperability Resources (FHIR) \cite{fhir_r4} is an HL7 standard for
exchanging healthcare data as modular \emph{resources} (for example \texttt{Patient},
\texttt{Condition}, \texttt{Observation}, \texttt{MedicationRequest}, and
\texttt{Procedure}). Resources reference one another, and clinical meaning is carried by
\texttt{CodeableConcept}/\texttt{Coding} elements, each pairing a code \texttt{system} URI
with a \texttt{code}. Crucially for this work, numeric results are first-class
\texttt{Observation} resources with a \texttt{valueQuantity}, so laboratory values and vital
signs are available as structured data. The release used throughout this thesis is FHIR~R4;
Section~\ref{sec:data-source} records how that was established from the exported instances.

Because the remainder of the thesis reasons directly over this structure, an abridged
example is given in Listing~\ref{lst:fhir}. It is taken verbatim from a generated bundle,
with unused elements elided.

\begin{lstlisting}[caption={Abridged FHIR R4 resources from a generated patient bundle:
a coded diagnosis and a numeric laboratory result.},label={lst:fhir},
basicstyle=\ttfamily\scriptsize,breaklines=true,frame=single,captionpos=b]
{ "resourceType": "Condition",
  "id": "22f599b6-...-8e66-1f64bdc5c863",
  "code": {                              // WHAT the finding is
    "coding": [ { "system":  "http://snomed.info/sct",
                  "code":    "414545008",
                  "display": "Ischemic heart disease (disorder)" } ],
    "text": "Ischemic heart disease (disorder)" },
  "subject": { "reference": "urn:uuid:22f599b6-...-e586-..." },   // WHOSE it is
  "onsetDateTime": "2013-05-24T09:14:31+07:00" }

{ "resourceType": "Observation",
  "status": "final",
  "code": {                              // WHICH measurement
    "coding": [ { "system":  "http://loinc.org",
                  "code":    "18262-6",
                  "display": "Cholesterol in LDL [Mass/volume] ..." } ] },
  "subject": { "reference": "urn:uuid:22f599b6-...-e586-..." },
  "valueQuantity": {                     // THE NUMBER -- machine-comparable
    "value": 91.35, "unit": "mg/dL",
    "system": "http://unitsofmeasure.org", "code": "mg/dL" } }
\end{lstlisting}

Three properties of this structure carry the design of the whole system. First, resources
are \emph{linked by reference}: a \texttt{Condition} does not contain its patient, it points
at one, which is why the graph model of Chapter~III turns references into edges. Second, clinical meaning is carried \emph{not by the field name but by the coding}. The same
\texttt{code} element holds a SNOMED~CT concept here and would hold an RxNorm or LOINC
concept elsewhere. That is why every coding is hoisted onto a shared node and bridged to a
common concept space. Third, and decisively for this thesis, a laboratory result is not a
sentence but a \texttt{valueQuantity} with a number, a unit, and a unit system. A flattened
tabular export carries those numbers too, but strips the structure they arrive in: a
multi-component reading is split across rows that must be rejoined before a threshold can be
evaluated. The FHIR resource keeps the value, its unit and its panel membership together,
which is what makes it possible to decide a recommendation that depends on a threshold rather
than merely on the presence of a diagnosis (Section~\ref{sec:flattening}).

\subsection{Clinical terminologies and the UMLS}
Different resource types are coded in different vocabularies: SNOMED~CT for clinical findings
and procedures \cite{snomedct}, RxNorm for medications \cite{rxnorm}, and LOINC for
laboratory and vital-sign observations \cite{loinc}. The Unified Medical Language System
(UMLS) Metathesaurus \cite{umls} integrates these and many other vocabularies and assigns
each concept a Concept Unique Identifier (CUI), which provides the common concept space into
which every EHR code and guideline target is mapped in this thesis.

\subsection{Knowledge graphs and property-graph databases}
A knowledge graph represents entities as nodes and their relationships as labelled edges. In
a property-graph database such as Neo4j \cite{neo4j}, nodes and edges may carry key--value
properties, and data are queried with the declarative Cypher language \cite{cypher} using
pattern matching. Property graphs are well suited to fusing heterogeneous sources because new
node and edge types can be added incrementally and multi-hop relationships can be traversed
directly.

\subsection{Clinical natural language processing: NER and entity linking}
Extracting structured information from unstructured biomedical text requires named-entity
recognition (NER), which identifies spans of clinical interest, and entity linking, which
maps those spans to a knowledge base such as the UMLS. The scispaCy toolkit \cite{scispacy}
provides biomedical NLP models and UMLS/MeSH entity linkers used in this work to process the
literature corpus and to assist terminology bridging; the model used here is
\texttt{en\_core\_sci\_md}, whose vectors are trained over OntoNotes~5, Common Crawl and
GENIA~1.0.

\subsection{Association strength: pointwise mutual information}
Pointwise mutual information (PMI) quantifies how much more often two events---here, a drug
and a condition---co-occur than would be expected if they were independent \cite{pmi}.
For two concepts $x$ and $y$ it is defined as

\begin{equation}
\mathit{PMI}(x, y) \;=\; \log \frac{P(x, y)}{P(x)\,P(y)} .
\label{eq:pmi}
\end{equation}

The probabilities are estimated by document frequency over the abstract corpus. Writing $N$
for the number of abstracts, $\mathit{df}(x)$ for the number mentioning concept $x$, and
$c(x,y)$ for the number mentioning both, the estimator substituted into
Equation~\ref{eq:pmi} is

\begin{equation}
\mathit{PMI}(x, y)
\;=\; \log_{10} \frac{c(x,y)/N}{\bigl(\mathit{df}(x)/N\bigr)\bigl(\mathit{df}(y)/N\bigr)}
\;=\; \log_{10} \frac{c(x,y)\;N}{\mathit{df}(x)\,\mathit{df}(y)} .
\label{eq:pmi-df}
\end{equation}

Base~10 is used, so a value of $1$ denotes a pair co-occurring ten times more often than
independence would predict and a value of $0$ denotes exact independence; negative values
denote co-occurrence rarer than chance. The right-hand form of Equation~\ref{eq:pmi-df} is
the one computed: two of the three factors of $N$ cancel, leaving a single $N$ in the
numerator, so the implementation materialises the corpus size once alongside the three
counts.

PMI is used to rank drug--condition associations mined from the literature. It is applied
with a minimum article-count floor: because $c(x,y)$ appears in the numerator, a pair
occurring in a single abstract can attain a high score purely by both concepts being rare,
so pairs below the floor are discarded rather than ranked. The floor is three articles: a
drug--condition pair must be co-mentioned in at least three distinct abstracts before it is
eligible for ranking. It is applied when the briefing is assembled rather than when the
co-occurrence edges are built, so the literature layer retains every pair it observed and
the floor can be varied without re-materialising the graph.

\subsection{Retrieval-augmented generation and grounded summarisation}
Retrieval-augmented generation (RAG) \cite{rag_lewis2020} conditions a language model's output
on retrieved context; classical RAG retrieves unstructured passages via dense (vector)
similarity, yielding \emph{approximate} context. This thesis instead uses a deterministic
knowledge-graph retriever, so the retrieved facts are exact and carry provenance, and the
language model performs \emph{grounded summarisation}---re-phrasing the facts without adding
information. Faithfulness denotes the property that every statement in the output is supported
by the retrieved facts; its violation is termed hallucination.

\subsection{Agreement measurement: Cohen's kappa}
Cohen's $\kappa$ \cite{cohen1960} measures the agreement between two raters (or between an
automated method and a gold standard) corrected for the agreement expected by chance. Values
above~0.8 are conventionally interpreted as near-perfect agreement, and values of 0.21--0.40
as fair agreement \cite{landis1977}. Kappa is
also affected by category prevalence and by bias between raters, so indicators of both should
be reported alongside it \cite{byrt1993}. It is used in this work to
validate the automated reference oracle against manual labels, and to measure agreement
between clinicians in the review of recommendation appropriateness.

\subsection{Synthetic patient data: Synthea}
Synthea \cite{synthea} is an open-source generator of synthetic but statistically realistic
patient records, produced by simulating disease progression through modular models and
exportable as HL7 FHIR bundles and CSV files. Because the data are synthetic, they carry no
protected health information and support fully reproducible research.

\section{Comparative Analysis of Methods}

The design choices of this thesis are justified by comparison with the alternatives. A purely
\emph{rule-based} CDSS can encode guideline logic but integrates poorly with heterogeneous
coded EHR data and is brittle to maintain; a \emph{knowledge-graph} approach with UMLS
bridging instead matches coded patient data to guideline concepts through a shared concept
space and traverses multi-source relationships directly. For the natural-language output, a
\emph{pure-LLM} summariser risks hallucination and a \emph{vector-RAG} summariser grounds only
on approximately retrieved passages; a \emph{deterministic knowledge-graph retrieval} followed
by LLM phrasing yields exact, provenance-bearing facts. Its output can therefore be checked
against a closed set of facts, which is essential in a clinical setting. Such a check does not
yet cover the semantic faithfulness of every sentence. Finally, evaluating on
\emph{FHIR-native} structured observations keeps each value bound to its unit, its unit system
and its panel, so an observation-dependent recommendation can be gated on what was measured
without first reassembling it from a flattened export---the mechanism by which precision is
improved. Table~\ref{tab:method-comparison} summarises this comparison.

\begin{table}[h]
\centering
\caption{Design alternatives for CCD decision support, compared on the criteria that
motivate this thesis. The relational row is included because it is the closest alternative
to the design adopted here, and is discussed separately below.}
\label{tab:method-comparison}
\footnotesize
\setlength{\tabcolsep}{3pt}
\begin{tabular}{p{2.0cm}p{2.28cm}p{2.28cm}p{2.28cm}p{2.28cm}p{2.28cm}}
\hline
\textbf{Approach} & \textbf{Coded EHR integration} & \textbf{Terminology bridging} &
\textbf{Numeric observation gating} & \textbf{Multi-source fusion} &
\textbf{Output faithfulness} \\
\hline
Rule-based CDSS / CIG &
Rare in practice: 2 of 15 evaluated systems \cite{papadopoulos2022} &
Hand-coded per rule &
Expressible, but hand-coded per rule &
Guideline only &
Not applicable (no NL output) \\
Relational EHR store with SQL &
Native &
Requires explicit mapping tables &
Native &
New schema and migration per source &
Not applicable (no NL output) \\
Vector-RAG over guideline text &
Text only; codes not used &
Embedding similarity &
Not represented &
Corpus-level, not entity-level &
Approximate; grounded on retrieved passages \\
Pure-LLM prompting &
Prompt text only &
Implicit and unverifiable &
Not represented &
None &
Unconstrained; hallucination risk \\
\textbf{This work} &
FHIR-native, structured &
UMLS shared concept space &
Gated on measured values &
EHR, guideline and literature in one graph &
Deterministic retrieval; automatically verified \\
\hline
\end{tabular}
\end{table}

One row of Table~\ref{tab:method-comparison} deserves separate comment, because it is the
alternative closest to the design adopted here. A relational store over the same FHIR data would also hold numeric observations natively.
The same canonical inference path could be expressed there as an equivalent multi-way join.
The choice of a property graph is therefore one of \emph{representation fit} rather than of
raw capability. The claim made here is not that the relational alternative is incapable. Three considerations nonetheless favour the graph.
First, FHIR defines well over a hundred resource types \cite{fhir_r4}. A faithful relational
schema needs a table per type and a migration for each newly ingested type. The dual-label
\texttt{:Resource:}\textit{Type} design of Chapter~III, by contrast, admits any resource
type through a single loader path. Second, the three fused sources---a patient record, a set
of conditional recommendations, and a literature co-occurrence network---share little
structure. A further layer can also be attached to a graph without disturbing the existing
ones. Third,
and most relevant to RQ3, a traversal written as a chain of clinical relationships is
auditable as a chain of clinical reasoning, which matters for a system whose output must carry
provenance. The contributions claimed in this thesis---the concept bridge, the observation
gates, and deterministic retrieval---rest above this choice rather than depending on it.

\section{Research Questions}

Each problem stated in Chapter~I is turned into a question that can be answered
empirically. The research questions below therefore correspond one-to-one to those
problems, and map to the research objectives:

\begin{enumerate}
    \item[RQ1.] How can a unified FHIR-native knowledge graph fusing EHR, guideline, and
    literature be constructed for CCD, with clinical terminologies bridged into a common
    concept space? (maps to RO1)
    \item[RQ2.] Does loading structured observations improve the recommendation system's
    guideline concordance relative to a condition/text-only baseline? (maps to RO2)
    \item[RQ3.] How faithful is the LLM-generated grounded briefing, and to what extent can
    deterministic knowledge-graph retrieval guarantee that faithfulness? (maps to RO3)
\end{enumerate}
